\section{Rejected Tails}
\label{sec:tails}

This section explains the proof of \cref{thm:tails}: a reader with a hole
has a hole whose coordinates with immortal endpoints have short rejected
tails. If $k=0$, the empty tuple is a hole without tails, and
\cref{thm:tails} holds with $\TailB(0,N)=0$. So let $k\ge1$. We fix a
structural reader $M$ with $k$ starts and $N$ states and a hole
$W=(w_1,\dots,w_k)$, with first rejected cuts $J_i$. The endpoint of $w_1$
is a rejected state, which differs from $\dead$, so $N\ge2$. Full proofs are
in \cref{app:tails,app:compression}.

\subsection{The Problem and the Plan}\label{sec:tails-plan}

Let $i$ be a coordinate with an immortal endpoint. After the cut $J_i$ its
run stays in rejected states, and no cut of the tail is the target of an
admitted pair. With one word this makes the tail useless, and it can be cut
off near $J_i$. With several words the tail still matters in the channels
where $i$ is the \emph{host}: the rest of another word $w_j$ after a cut
$b\le J_j$ may be a prefix of the rest of $w_i$ after a cut $a$ with an
admitted pair of states, unless $w_i$ ends too early or its text deviates.
The tail of $w_i$ may be exactly what makes the texts deviate. Cutting it
off, or deleting a loop from it, can therefore create an admitted pair, and
how long the tail must be can depend on the other words, whose lengths are
not bounded at this stage.

The plan is to make the tails \emph{periodic} first, and to shorten them
afterwards.
\begin{enumerate}
\item \emph{Rays.} We replace the tail of $w_i$, from a cut $c_i\ge J_i$ on,
  by a periodic continuation $\tau_i^\omega$, where $\tau_i$ is a factor of
  $w_i$ that is a loop of the run. The infinite word
  $\wpre{w_i}{c_i}\tau_i^\omega$ is a \emph{ray}
  (\cref{sec:tails-rays}).
\item \emph{No new comparisons with the other words.} The factor $\tau_i$ is
  chosen so that repeating it creates no comparison with the words that are
  not replaced (\cref{sec:tails-closed}). Which words are replaced is
  decided by measuring how far each tail is from being periodic, at a
  suitable scale (\cref{sec:tails-levels}).
\item \emph{Compatible rays.} Two rays can agree forever from some cuts on.
  Whether the corresponding pair of states is compared then depends only on
  the lengths at which the rays are cut off, through a difference constraint.
  We choose the rays so that these constraints have a solution
  (\cref{sec:tails-choice}). This yields a hole whose immortal tails consist
  of a bounded number of periodic stretches and a bounded number of other
  letters.
\item \emph{Compression.} Long periodic stretches are shortened by deleting
  blocks of periods, in all stretches with the same period pattern at once
  (\cref{sec:tails-compression}).
\end{enumerate}

\subsection{Rays and Difference Constraints}\label{sec:tails-rays}

\begin{definition}[Rays]\label{def:ray}
A \emph{ray} is an ultimately periodic infinite word $\Omega=u\tau^\omega$,
with $\tau$ nonempty. Its \emph{period} $\per(\Omega)$ is the least $p\ge1$
such that $\Omega[x+p]=\Omega[x]$ for all sufficiently large positions $x$,
and its \emph{onset} $\onset(\Omega)$ is the least position from which this
holds. We write $\wpre\Omega L$ for its prefix of length $L$ and
$\wsuf\Omega a$ for its infinite rest after the cut $a$.
\end{definition}

For example, $00(10)^\omega=0\,(01)^\omega$ is a ray with period $2$ and
onset $1$. The presentation $u\tau^\omega$ of a ray is not unique, but its
period and its onset are.

A ray is read by the reader like a word, from the start state of its
coordinate; we consider rays whose runs never reach $\dead$. Two rays that
agree on a long stretch agree forever:

\begin{lemma}[Comparisons of rays]\label{lem:ray-compare}
Let $\Omega$ and $\Omega'$ be rays and $a,b$ cuts.
\begin{enumerate}
\item\label{it:ray-fw} If $\wsuf\Omega a$ and $\wsuf{\Omega'}b$ agree on
  their first $\max(\onset(\Omega),\onset(\Omega'))+\per(\Omega)+
  \per(\Omega')$ letters, then $\wsuf\Omega a=\wsuf{\Omega'}b$.
\item\label{it:ray-onset} If $\wsuf\Omega a=\wsuf{\Omega'}b$ and
  $b<\onset(\Omega')$, then $\onset(\Omega)-a=\onset(\Omega')-b$ and
  $\per(\Omega)=\per(\Omega')$.
\end{enumerate}
\end{lemma}

Now let rays $\Omega_i$ be given for the coordinates $i$ of a set $I$, each
with the first rejected cut $J_i$, and consider the prefixes
$\wpre{\Omega_i}{L_i}$ for lengths $L_i$ that exceed $J_i$ by more than the
bound of \cref{lem:ray-compare}\eqref{it:ray-fw}. An admitted pair on the
channel from $i$ to $j$ has a target cut $b\le J_j$, so the rest of the
target is long and the two rays agree forever from the two cuts. Conversely,
if they agree forever, the prefixes are compared at these cuts exactly when
the rest of the target is not longer than the rest of the host
(\cref{fig:ray}).

\begin{definition}[Ray system]\label{def:ray-system}
The \emph{ray system} of $(\Omega_i)_{i\in I}$ has a constraint
\[
  L_j\;\ge\;L_i+(b-a+1)
\]
for all $i\ne j$ in $I$ and all cuts $a$ of $\Omega_i$ and $b\le J_j$ of
$\Omega_j$ such that the pair of states at $a$ and $b$ is admitted and
$\wsuf{\Omega_i}a=\wsuf{\Omega_j}b$. A \emph{self event} of $\Omega_i$ is a
pair of cuts $a\le b\le J_i$ with an admitted pair of states and
$\wsuf{\Omega_i}a=\wsuf{\Omega_i}b$.
\end{definition}

The states along a ray are ultimately periodic, so the constraints whose host
cut $a$ is smaller than $\onset(\Omega_i)+N\per(\Omega_i)$ imply all others,
and the target cuts are at most $J_j$ (\cref{app:tails-rays}). So the ray
system is equivalent to a finite system of difference constraints, whose
graph has an edge $i\to j$ of weight $b-a+1$ for each constraint. For
example, let $\Omega_1=(01)^\omega$ and $\Omega_2=1(01)^\omega$, and suppose
that the states at the cut $a=1$ of $\Omega_1$ and the cut $b=0$ of
$\Omega_2$ form an admitted pair. The rests after these cuts are both
$(10)^\omega$. In the prefixes of lengths $L_1$ and $L_2$, the rest of the
target has $L_2$ letters and the rest of the host $L_1-1$, so $(1,0)$ is a
comparison exactly when $L_2\le L_1-1$. It is avoided by the constraint
$L_2\ge L_1$, which is $L_2\ge L_1+(b-a+1)$: the target must outlast the
host.

\begin{lemma}[Prefixes of rays]\label{lem:ray-prefixes}
\begin{enumerate}
\item\label{it:rp-hole} If no ray has a self event and the lengths $L_i$
  satisfy the ray system and exceed $J_i+\max_j\onset(\Omega_j)+2\max_j
  \per(\Omega_j)$, with the maxima over $j\in I$ (a maximum over the empty
  set is zero, and for $I=\emptyset$ the assertion is vacuous), then no
  channel between two coordinates of $I$, equal or
  not, of the tuple $(\wpre{\Omega_i}{L_i})_{i\in I}$ contains an admitted
  pair.
\item\label{it:rp-solve} A system of difference constraints has a solution
  in natural numbers above any given lower bounds if and only if its graph
  has no cycle of positive weight.
\end{enumerate}
\end{lemma}

%% ray -- a ray of a coordinate, and a constraint between two rays.
\begin{figure}[t]
\centering
\begin{tikzpicture}[font=\scriptsize, >={Stealth[length=1.8mm]}]
\def\h{0.34}
%% (a) a ray
\begin{scope}
  \draw (0,0) rectangle (3.0,\h);
  \node at (1.0,0.5*\h) {head};
  \draw (2.0,0) -- (2.0,\h);
  \foreach \x in {3.0,4.0,5.0,6.0} {
    \draw[fill=black!15] (\x,0) rectangle (\x+1.0,\h);
    \node at (\x+0.5,0.5*\h) {$\tau$};
  }
  \node at (7.35,0.5*\h) {$\cdots$};
  \draw (2.0,-0.08) -- (2.0,-0.3) node[below] {$J_i$};
  \draw (3.0,-0.08) -- (3.0,-0.3) node[below] {$c_i$};
  \draw (4.0,-0.08) -- (4.0,-0.3) node[below] {$c_i'$};
  \draw[decorate, decoration={brace, amplitude=4pt}] (0,\h+0.08) -- (4.0,\h+0.08)
    node[midway, above=4pt] {the word $w_i$ up to $c_i'$};
  \draw (6.5,-0.08) -- (6.5,-0.3) node[below] {$L_i$};
  \draw[densely dashed] (6.5,-0.05) -- (6.5,\h+0.05);
  \node[left] at (0,0.5*\h) {$\Omega_i$};
  \node at (3.6,-1.05) {(a) the ray $\Omega_i=\wpre{w_i}{c_i}\,\tau^\omega$ of a factor
    $\tau=\wsub{w_i}{c_i}{c_i'}$ and its prefix of length $L_i$};
\end{scope}
%% (b) a constraint between two rays
\begin{scope}[yshift=-3.0cm]
  % host ray
  \draw (0,0) rectangle (7.2,\h);
  \fill[black!15] (1.4,0) rectangle (7.2,\h);
  \draw (1.4,0) -- (1.4,\h);
  \node[left] at (0,0.5*\h) {$\Omega_i$};
  \draw (1.4,\h+0.08) -- (1.4,\h+0.3) node[above] {$a$};
  \draw[densely dashed] (5.3,-0.05) -- (5.3,\h+0.05);
  \draw (5.3,\h+0.08) -- (5.3,\h+0.3) node[above] {$L_i$};
  % target ray
  \draw (1.0,-1.1) rectangle (7.2,-1.1+\h);
  \fill[black!15] (1.4,-1.1) rectangle (7.2,-1.1+\h);
  \draw (1.4,-1.1) -- (1.4,-1.1+\h);
  \node[left] at (1.0,-1.1+0.5*\h) {$\Omega_j$};
  \draw (1.4,-1.18) -- (1.4,-1.4) node[below] {$b\le J_j$};
  \draw[densely dashed] (6.2,-1.15) -- (6.2,-1.1+\h+0.05);
  \draw (6.2,-1.18) -- (6.2,-1.4) node[below] {$L_j$};
  \node at (4.3,-0.38) {$\wsuf{\Omega_i}{a}=\wsuf{\Omega_j}{b}$};
  \node at (7.45,0.5*\h) {$\cdots$};
  \node at (7.45,-1.1+0.5*\h) {$\cdots$};
  \node at (3.6,-2.15) {(b) an admitted pair at $(a,b)$ with equal infinite rests is harmless
    if $L_j-b>L_i-a$};
\end{scope}
\end{tikzpicture}
\caption{(a) A ray continues a word periodically after its first rejected
cut. (b) If the states at the cuts $a$ and $b$ are admitted and the rays
agree forever from these cuts, the prefixes of lengths $L_i$ and $L_j$ carry
this pair on the channel from $i$ to $j$ unless the rest of the target is
longer than the rest of the host. This gives the difference constraint
$L_j\ge L_i+(b-a+1)$.}
\label{fig:ray}
\end{figure}


So rays turn the unbounded problem of avoiding each other into a finite
graph condition. The remaining difficulty is to choose the rays so that
there is no positive cycle, and so that they create no comparisons with the
words that stay as they are.

\subsection{Closed Factors}\label{sec:tails-closed}

Let $x$ be the word of an eligible coordinate, with the first rejected cut
$J$, and let $t_1,\dots,t_m$ be words, the \emph{target words}: they stand
for the words of the other coordinates, whose rests may be compared with
rests of $x$. We want cuts $J\le c<c'$ of $x$ such that replacing the tail
of $x$ after $c$ by repetitions of $\tau=\wsub xc{c'}$ creates no comparison
in which this coordinate is the host and the target is one of the target
words. A comparison ends where the target word ends, so what matters at a
cut $e$ of the host is which rests of target words end there, and with which
state before them. For rests of bounded length this is finite information.

\begin{definition}[Suffix nodes and closed factors]\label{def:node}
Let $B\ge0$ be a \emph{window}. For a target word $t$, $v_B(t)$ denotes its
suffix of length $B$ if $\len t\ge B$; if $\len t<B$, there is no such word,
and we say that $v_B(t)$ has no occurrences. The \emph{suffix dictionary}
$\Dict_B$ consists of the empty word and the factors of the words $v_B(t_j)$,
or of $t_j$ itself if $\len{t_j}<B$. The \emph{suffix node} of a cut $e$ of
$x$ is the pair
\[
  \node_B(e)=\bigl(\sigma,\;q(e-\len\sigma)\bigr),
  \qquad\sigma\text{ the longest suffix of }\wpre xe\text{ in }\Dict_B,
\]
where $q(\cdot)$ is the state of the run of $x$. A factor $\wsub xc{c'}$
with $c<c'$ is \emph{closed at $B$} if $\node_B(c)=\node_B(c')$, and it is
\emph{free at $B$} if no occurrence of a word $v_B(t_j)$ ends at a cut of
$(c,c']$.
\end{definition}

\begin{example}\label{ex:nodes}
Let $t=0110$ be the only target word and $B=2$. Then $v_2(t)=10$ and
$\Dict_2=\{\eps,1,0,10\}$. In the host word $x=101100$, the longest suffix
of $\wpre x2=10$ in the dictionary is $10$, so $\node_2(2)=(10,q(0))$; at
the cut $3$ it is $1$, and $\node_2(3)=(1,q(2))$; at the cut $5$ it is again
$10$, and $\node_2(5)=(10,q(3))$. The factor $\wsub x25=110$ is closed at
$2$ if $q(0)=q(3)$. It is not free at $2$, because an occurrence of
$v_2(t)=10$ ends at the cut $5$.
\end{example}

The dictionary has at most $mB(B+1)/2+1$ words, whatever the alphabet, so
there are at most $N(mB(B+1)/2+1)$ suffix nodes. Since the dictionary is
closed under factors, the suffix node after one more letter is a function
of the suffix node before and of the letter, as in the automaton of Aho and
Corasick \cite{AhoCorasick75}. Hence along a repetition of a closed factor
the suffix nodes repeat (\cref{fig:closed-factor}).

\begin{lemma}[Pumping a closed factor]\label{lem:pump}
Let $\wsub xc{c'}$ be closed and free at $B$, with $c\ge J$, and let
$x'=\wpre xc\,\wpre{(\tau^\omega)}L$ with $\tau=\wsub xc{c'}$ and $L\ge
c'-c$. Then $x'$ is eligible with the first rejected cut $J$ and an immortal
endpoint, and for every target word $t_j$ and all cuts $b$ of $t_j$ and $a$ of
$x'$: if $\wsuf{t_j}b$ is a prefix of $\wsuf{x'}a$, then it is a prefix of
$\wsuf x{a^\circ}$ for a cut $a^\circ$ of $x$ that carries the same state as
the cut $a$ of $x'$.
\end{lemma}

%% closed-factor -- suffix nodes, a closed factor free of target occurrences.
\begin{figure}[t]
\centering
\begin{tikzpicture}[font=\scriptsize, >={Stealth[length=1.8mm]}]
\def\h{0.34}
% host word
\draw (0,0) rectangle (10.0,\h);
\node[left] at (0,0.5*\h) {$x$};
\draw[fill=black!15] (4.4,0) rectangle (7.6,\h);
\node at (6.0,0.5*\h) {$\tau$};
% suffix in dictionary at c and at c'
\draw[pattern=north east lines, pattern color=black!45] (3.3,0) rectangle (4.4,\h);
\draw[pattern=north east lines, pattern color=black!45] (6.5,0) rectangle (7.6,\h);
\node[above] at (3.85,\h) {$\sigma$};
\node[above] at (7.05,\h) {$\sigma$};
\draw (4.4,-0.08) -- (4.4,-0.3) node[below] {$c$};
\draw (7.6,-0.08) -- (7.6,-0.3) node[below] {$c'$};
\node[circle, draw, inner sep=1pt] (q1) at (3.3,-0.95) {$q$};
\node[circle, draw, inner sep=1pt] (q2) at (6.5,-0.95) {$q$};
\draw[->] (q1) -- (3.3,-0.08);
\draw[->] (q2) -- (6.5,-0.08);
\node[anchor=west] at (7.85,-0.95) {$\node_B(c)=\node_B(c')=(\sigma,q)$};
% a target suffix ending at an old cut
\draw (0.9,1.25) rectangle (2.6,1.25+\h);
\node at (1.75,1.25+0.5*\h) {$v_B(t)$};
\draw[densely dotted] (0.9,1.25) -- (0.9,\h);
\draw[densely dotted] (2.6,1.25) -- (2.6,\h);
\node[anchor=west] at (2.75,1.25+0.5*\h) {an occurrence of the suffix of length $B$ of a target};
\draw[decorate, decoration={brace, mirror, amplitude=4pt}] (4.45,-1.5) -- (7.6,-1.5)
  node[midway, below=4pt] {no such occurrence ends in $(c,c']$};
\end{tikzpicture}
\caption{A closed factor. The cuts $c<c'$ of the host word $x$ have the same
suffix node: the same longest suffix $\sigma$ in the suffix dictionary of the
targets at the window $B$, and the same state $q$ before it. If moreover no
occurrence of the suffix of length $B$ of a target ends in $(c,c']$, the
factor $\tau=\wsub xc{c'}$ can be repeated any number of times without
creating a comparison with a target (\cref{lem:pump}).}
\label{fig:closed-factor}
\end{figure}


A rest of a target word of length at most $B$ that ends at a cut of $x'$ lies in
the dictionary, so it is a suffix of the word of the suffix node, and it ends
at a cut of $x$ with the same node. A longer rest ends with $v_B(t_j)$, which
by freedom ends only at cuts up to $c$, where $x'$ and $x$ coincide.

Closed factors exist by the pigeonhole principle; the point is to find one
that is free. The occurrences of $v_B(t)$ can be avoided if they are
sparse, and they are sparse unless $t$ has a periodic suffix at the scale
of $B$. This is measured by \emph{suffix chains}: for a window $b$ with
$1\le b\le\len t$ let $\gap_t(b)$ be the distance from the end of $t$ to the
end of the previous occurrence of $v_b(t)$ in $t$, if there is one; a chain
is a sequence $b_0<b_1<\cdots$ of windows with $b_{r+1}=b_r+\gap_t(b_r)$,
and its number of \emph{steps} is the number of its members after the first.
Along a chain the suffixes $v_{b_r}(t)$ are prefixes of one another.

\begin{example}\label{ex:chain}
Let $t=ababab$. The suffix $v_1(t)=b$ also ends two letters before the end
of $t$, so $\gap_t(1)=2$ and the chain from the window $1$ continues with
$3$. The suffix $v_3(t)=bab$ again ends two letters before the end, so
$\gap_t(3)=2$, and the next window is $5$. The suffix $v_5(t)=babab$ has no
other occurrence, and the chain $1<3<5$ stops after two steps
(\cref{fig:chain}). A periodic suffix thus gives a chain whose gap is the
period, with about as many steps as the suffix has periods. In contrast,
for $t=aabbab$ the chain from the window $1$ is $1<3$: the gap of $b$ is
again $2$, but $bab$ occurs only once. A chain stops at a window whose
suffix occurs only at the end of $t$, and a long chain with small gaps
indicates a long periodic suffix.
\end{example}

%% chain -- the suffix chain of ababab from the window 1.
\begin{figure}[t]
\centering
\begin{tikzpicture}[font=\scriptsize, x=0.62cm, y=1cm]
\def\h{0.36}
\foreach \c [count=\i from 0] in {a,b,a,b,a,b} {
  \draw (\i,0) rectangle (\i+1,\h);
  \node at (\i+0.5,0.5*\h) {$\c$};
}
\node[left] at (0,0.5*\h) {$t$};
\foreach \i in {0,...,6} { \node[above, font=\tiny] at (\i,\h) {$\i$}; }
%% rows: window b, suffix [6-b,6], previous occurrence [6-b-2,6-2]
\foreach \b/\y in {1/-0.45,3/-1.2} {
  \pgfmathsetmacro{\l}{6-\b}
  \draw[line width=1.6pt] (\l,\y) -- (6,\y);
  \draw[line width=1.6pt, black!45] (\l-2,\y-0.22) -- (4,\y-0.22);
  \draw[densely dotted] (6,\y) -- (6,0);
  \draw[densely dotted] (4,\y-0.22) -- (4,0);
  \draw[<->, >={Stealth[length=1.4mm]}] (4,\y-0.42) -- node[below, inner sep=1.5pt] {gap $2$} (6,\y-0.42);
}
\draw[line width=1.6pt] (1,-2.05) -- (6,-2.05);
\node[anchor=west] at (6.6,-0.5) {window $1$: the suffix $b$ and its previous occurrence};
\node[anchor=west] at (6.6,-1.25) {window $3=1+2$: the suffix $bab$ and its previous occurrence};
\node[anchor=west] at (6.6,-2.05) {window $5=3+2$: the suffix $babab$ occurs only once};
\end{tikzpicture}
\caption{The suffix chain $1<3<5$ of $t=ababab$ (\cref{ex:chain}). At each
window the suffix (black) is compared with its previous occurrence (gray);
the distance between their ends is the gap, which is added to the window.
The suffixes $b$, $bab$, $babab$ are prefixes of one another.}
\label{fig:chain}
\end{figure}


\begin{lemma}[Closed factors free of the target words]\label{lem:closed-free}
For every number $s$ there are computable functions $\Lb_m$ and $\Wb_m$,
which depend on $s$ and are nondecreasing in the window, with the following
property. Let $1\le b_0\le b$, and suppose that every suffix chain of every
target word that starts at a window of at least $b_0$ has at most $s$
steps. Then every sequence of $\Lb_m(b)$
consecutive cuts of $x$ contains cuts $c<c'$ and there is a window $B$ with
$b\le B\le\Wb_m(b)$ such that $\wsub xc{c'}$ is closed and free at $B$ and
$B<c'-c\le\Fb(B)$, where $\Fb(B)=(B+1)\,N\,(mB(B+1)/2+1)$.
\end{lemma}

The proof treats the target words one after the other. For one target word
it follows a suffix chain as long as the gaps stay below the length needed
for the remaining ones; the chain stops within $s$ steps, at a window where the
occurrences of the suffix are far apart. Either a long stretch of the given
cuts is free at this window, or the occurrences come so regularly that a
slightly larger window has no occurrence at all after the first one. The
base case is a pigeonhole on the suffix nodes of cuts that are $B+1$ apart,
which makes the closed factor longer than its window. The functions $\Lb_m$
are polynomials in the window whose degree depends on $s$ and $m$; this
matters for the complexity bound (\cref{sec:complexity}).

\subsection{Scales of Periodicity}\label{sec:tails-levels}

%% types -- the three types of coordinates at a stable scale.
\begin{figure}[t]
\centering
\begin{tikzpicture}[font=\scriptsize, >={Stealth[length=1.8mm]}, xscale=0.86]
\def\h{0.32}
\tikzset{pf/.style={fill=black!18}, ex/.style={fill=black!75}}
%% fixed
\begin{scope}
  \draw (0,0) rectangle (7.4,\h);
  \node at (1.0,0.5*\h) {head};
  \draw (2.0,-0.08) -- (2.0,-0.28) node[below] {$J_i$};
  \draw (2.0,0) -- (2.0,\h);
  \draw[pf] (2.15,0) rectangle (4.2,\h);
  \draw[ex] (4.2,0) rectangle (4.3,\h);
  \draw[pf] (4.3,0) rectangle (6.3,\h);
  \draw[ex] (6.3,0) rectangle (6.45,\h);
  \draw[pf] (6.45,0) rectangle (7.4,\h);
  \node[left] at (0,0.5*\h) {fixed};
  \node[anchor=west] at (7.6,0.5*\h) {kept; a target word for the others};
\end{scope}
%% terminal
\begin{scope}[yshift=-1.15cm]
  \draw (0,0) rectangle (7.4,\h);
  \node at (1.0,0.5*\h) {head};
  \draw (2.0,-0.08) -- (2.0,-0.28) node[below] {$J_i$};
  \draw (2.0,0) -- (2.0,\h);
  \draw[pf] (2.2,0) rectangle (3.4,\h);
  \draw[ex] (3.4,0) rectangle (3.5,\h);
  \draw[pf] (3.5,0) rectangle (7.4,\h);
  \draw[pf, densely dashed] (7.4,0) rectangle (9.3,\h);
  \node at (9.6,0.5*\h) {$\cdots$};
  \draw[decorate, decoration={brace, mirror, amplitude=3pt}] (3.5,-0.06) -- (7.4,-0.06)
    node[midway, below=3pt] {long periodic suffix};
  \node[left] at (0,0.5*\h) {terminal};
  \node[anchor=west] at (9.9,0.5*\h) {its ray};
\end{scope}
%% escaping
\begin{scope}[yshift=-2.55cm]
  \draw (0,0) rectangle (7.4,\h);
  \node at (1.0,0.5*\h) {head};
  \draw (2.0,-0.08) -- (2.0,-0.28) node[below] {$J_i$};
  \draw (2.0,0) -- (2.0,\h);
  \draw[pattern=north east lines, pattern color=black!45] (2.0,0) rectangle (7.4,\h);
  \draw[pf] (4.4,0) rectangle (5.3,\h);
  \node at (4.85,0.5*\h) {$\tau$};
  \draw (4.4,-0.08) -- (4.4,-0.28) node[below] {$c_i$};
  \draw (5.3,-0.08) -- (5.3,-0.28) node[below] {$c_i'$};
  \node[left] at (0,0.5*\h) {escaping};
  %% the ray
  \draw (0,-1.0) rectangle (4.4,-1.0+\h);
  \draw[pattern=north east lines, pattern color=black!45] (2.0,-1.0) rectangle (4.4,-1.0+\h);
  \draw (2.0,-1.0) -- (2.0,-1.0+\h);
  \foreach \x in {4.4,5.3,6.2,7.1,8.0} {
    \draw[pf] (\x,-1.0) rectangle (\x+0.9,-1.0+\h);
    \node at (\x+0.45,-1.0+0.5*\h) {$\tau$};
  }
  \node at (9.2,-1.0+0.5*\h) {$\cdots$};
  \node[anchor=west] at (9.9,-1.0+0.5*\h) {its ray};
  \node[anchor=west] at (7.6,0.5*\h) {cut at a closed and free factor $\tau$};
\end{scope}
\end{tikzpicture}
\caption{The three types of coordinates at a stable scale. Periodic factors
with short periods are gray, exceptional positions black, and text without
such a structure is hatched. A fixed coordinate with an immortal endpoint
has a tail of few periodic factors and no long periodic suffix; its word is
kept. A terminal coordinate ends with a long periodic suffix and is
continued by the ray of this suffix. An escaping coordinate is far from
periodic at this scale; its tail is cut at one of many closed and free
factors $\tau$ and continued by $\tau^\omega$.}
\label{fig:types}
\end{figure}


Which coordinates should be replaced by rays, and at which scale? A tail
that already consists of few periodic stretches should be kept: its word is
a harmless target word, with short suffix chains. A tail that is far from
periodic offers many closed factors with different onsets to choose from.
The two cases must be separated by a gap in scale, because the closed
factors of one word have to avoid the suffixes of the others. We therefore
measure periodicity at a sequence of \emph{scales} $l=0,1,2,\dots$, each
much coarser than the one before. Scale $l$ comes with four numbers, defined
by a recursion on $l$ in \cref{app:tails-levels} and nondecreasing in $l$:
\begin{itemize}
\item $e_l$ bounds the periods that count as short at scale $l$;
\item $T_l$ bounds the number of positions of a tail that may lie outside
  its periodic factors;
\item $Y_l$ is the length from which a periodic suffix counts as long;
\item $b_l$ is the \emph{initial window} of scale $l$, the smallest window
  at which closed factors are searched.
\end{itemize}
The recursion makes each scale dominate the previous ones: for instance
$b_{l+1}>(N+1)\,e_l$, and $Y_{l+1}$ exceeds every length that can be
assembled from $T_l$ exceptional positions and boundedly many periods
$e_{l+1}$ (inequalities (L1)--(L5) of \cref{app:tails-levels}). The number
of periodic factors that we allow is governed by the constant
\[
  \Qc=4+2(k-1)(N-1).
\]
Its value is chosen for the count in \cref{sec:tails-choice}: a coordinate that
is not almost periodic offers at least $\Qc-1$ suitable rays, and each of
the other $k-1$ coordinates can spoil at most $2(N-1)$ of them, so one
remains because $\Qc-1>2(k-1)(N-1)$.

\begin{definition}[Almost periodic tails]\label{def:cover}
An eligible coordinate $i$ has property $F_l$ if all positions of its
rejected tail, except at most $T_l$, lie in at most $\Qc+1$ factors of $w_i$
with periods at most $e_l$. It has property $K_l$ if it has $F_l$ and the
longest suffix of $w_i$ with a period of at most $e_l$ is shorter than
$Y_l$.
\end{definition}

So $F_l$ says that the tail consists of few periodic stretches, up to few
exceptional positions, and $K_l$ says in addition that the word does not
end with a long periodic stretch. For instance, for three different
letters $a,b,c$, a tail $(ab)^{r}c\,(ba)^{r'}$ has $F_l$ for every $l$, with
two factors of period $2$ and one exceptional position, whatever $r$ and
$r'$ are, and it has $K_l$ exactly when $2r'<Y_l$. In contrast, a tail
fails $F_l$ when its periodic stretches with periods up to $e_l$ are so
many and so short that any $\Qc+1$ of them leave more than $T_l$ positions
uncovered.

Both properties are inherited from a scale to the next
(\cref{lem:levels-mono}); for $K_l$ this uses the periodicity lemma of Fine
and Wilf \cite{FineWilf65}. There are $2k$ properties $F_l,K_l$ of the $k$
coordinates, each of which changes at most once as $l$ grows, so among the
scales up to $3(2k+1)$ there is a \emph{stable} scale $l\ge2$: no coordinate
changes its status with respect to $F$ between $l-1$ and $l$, or with
respect to $K$ between $l-2$ and $l-1$. At a stable scale every coordinate
is of exactly one of three types (\cref{fig:types}).
\begin{itemize}
\item \emph{Fixed}: the endpoint is mortal, or the coordinate has $K_{l-2}$.
  Its word is kept. Its suffix chains from the initial window $b_l$ of
  scale $l$ on have at most $\Qc+N$ steps (\cref{lem:chains}); for a mortal
  endpoint this follows from \cref{lem:distinct}.
\item \emph{Terminal}: the endpoint is immortal and the coordinate has
  $F_{l-1}$ but not $K_{l-1}$. Its word ends with a long suffix of period at
  most $e_{l-1}$, and it is continued by the ray of this suffix. The ray
  contains the whole word $w_i$.
\item \emph{Escaping}: the endpoint is immortal and the coordinate does not
  have $F_l$. Its tail is cut at a suitable closed factor and continued
  periodically.
\end{itemize}
The number of steps of a suffix chain of a fixed word does not grow with the
scale; this is what keeps the recursion of the numbers $e_l,T_l,Y_l,b_l$,
and with it the final bound, elementary.

\subsection{Choosing Compatible Rays}\label{sec:tails-choice}

We fix a stable scale $l$. The target words are the words of the fixed
coordinates. For a terminal
coordinate, the last $pN!$ letters of its periodic suffix, $p$ the period,
form a factor that is closed and free at a suitable window
(\cref{lem:terminal}). For an escaping coordinate $v$, \cref{lem:closed-free}
provides a closed and free factor in every stretch of $\Lb_k(b_l)$ cuts of
the tail, a \emph{candidate}. If the rays of the candidates had at most $\Qc+1$ different
onsets, the periodic stretches from these onsets would cover the tail up to
$T_l$ positions, and $v$ would have $F_l$. So there are candidates whose
rays have at least $\Qc+2$ different onsets, of which at most three lie at
or before $J_v$ (\cref{lem:candidates}); this leaves $\Qc-1$ candidate rays
with different onsets after $J_v$.

It remains to exclude positive cycles in the ray system. A constraint $i\to
j$ requires the rays to agree forever from cuts $a$ and $b\le J_j$. If the
onset of $\Omega_j$ lies after $J_j$, \cref{lem:ray-compare} aligns the
onsets, and the stretch $X_j=\wsub{\Omega_j}{J_j}{\onset(\Omega_j)}$ of the
target between its first rejected cut and its onset is copied into the host
just before the onset of the host. Following a cycle backwards from an
escaping coordinate $v$, these stretches grow strictly, so the cycle cannot
return to $v$, unless one of three \emph{forbidden overlaps} of $v$ with
the ray of another coordinate $u$ occurs (\cref{def:overlaps}). In the
first, $X_v$ is a suffix of the text of the ray of $u$ before its onset,
the onset of $u$ lies after $J_u$, and $X_u$ is not longer than $X_v$. In
the second, $X_v$ is such a suffix, the onset of $u$ lies at or before
$J_u$, and the two rays have the same periodic part. In the third, the
roles are exchanged: $X_u$ is a nonempty suffix of the text of the ray of
$v$ before its onset, and $X_v$ is not longer than $X_u$.

\begin{lemma}[Compatible candidates]\label{lem:compatible}
\begin{enumerate}
\item\label{it:compat-count} For a terminal ray or an escaping ray chosen
  earlier, of a coordinate $u\ne v$, at most $2(N-1)$ candidates of $v$ with
  pairwise different onsets after $J_v$ have a forbidden overlap with it.
\item\label{it:compat-cycle} A ray with its onset after its first rejected
  cut that has no forbidden overlap with any other ray lies on no cycle of
  the ray system.
\end{enumerate}
\end{lemma}

The count in \eqref{it:compat-count} is again an application of
\cref{lem:distinct}: different candidates with the same kind of overlap give
cuts at or before a first rejected cut that carry different states. Since
$\Qc-1>2(k-1)(N-1)$, the escaping coordinates can be treated one after the
other, each taking one of its $\Qc-1$ candidates that has no forbidden
overlap with the at most $k-1$ rays chosen before. Cycles then pass through terminal rays only, and for these the
lengths of the original words satisfy all constraints, because $W$ is a hole.
So there is no positive cycle. Solving the ray system above sufficiently
large lower bounds and applying \cref{lem:pump,lem:ray-prefixes} gives:

\begin{proposition}[Holes with periodic tails]\label{prop:run-cover}
There are computable numbers $D$ and $E_0$, depending on $k$ and $N$, such
that every structural reader with a hole has a hole in which the rejected
tail of every coordinate with an immortal endpoint is covered, up to $E_0$
positions in total, by at most $\Qc+3$ maximal periodic factors with periods
at most $D$ and lengths at least $2D$.
\end{proposition}

\subsection{Compressing the Pieces}\label{sec:tails-compression}

A \emph{piece} of period $d$ in a tuple is a maximal factor of one word that
has period $d$, length at least $2d$ and a primitive \emph{root}, its first
$d$ letters; a word is primitive if it is not a power of a shorter word.
Pieces are the runs, or maximal repetitions, of combinatorics on words
\cite{KolpakovKucherov99}; we say piece because run already means the run of
an automaton. The maximal periodic factors of \cref{prop:run-cover} are
pieces. The pieces whose roots are conjugates (rotations) of a primitive
word $\zeta$ form the \emph{class} of $\zeta$. The \emph{follower} of a
piece is the rest of its word after it. For example, the word $abababb$ has
the piece $ababab$ of period $2$, with the root $ab$ and the follower $b$,
and the piece $bb$ of period $1$, with the root $b$ and the empty follower;
in $bababa$ the whole word is a piece of the class of $ab$, with the root
$ba$.

%% compression -- an admissible deletion on two pieces of a class with comparable followers.
\begin{figure}[!b]
\centering
\begin{tikzpicture}[font=\scriptsize, >={Stealth[length=1.8mm]}]
\def\h{0.34}
%% before
\begin{scope}
  % word i with a long piece P = [0.8,9.2], site [1.2,5.1]
  \draw (0,0) rectangle (11.0,\h);
  \node[left] at (0,0.5*\h) {$w_i$};
  \draw[fill=black!10] (0.8,0) rectangle (9.2,\h);
  \draw[fill=black!25] (1.2,0) rectangle (5.1,\h);
  \node at (3.15,0.5*\h) {deleted: $x_P$ blocks};
  \node at (7.1,0.5*\h) {piece $P$};
  \node at (10.1,0.5*\h) {follower};
  \draw[decorate, decoration={brace, amplitude=3pt}] (0.8,\h+0.06) -- (1.2,\h+0.06)
    node[midway, above=3pt] {$\beta_d$};
  % word j with a piece R = [5.3,9.2], site [5.7,6.6], comparable follower
  \draw (4.5,-1.3) rectangle (10.4,-1.3+\h);
  \node[left] at (4.5,-1.3+0.5*\h) {$w_j$};
  \draw[fill=black!10] (5.3,-1.3) rectangle (9.2,-1.3+\h);
  \draw[fill=black!25] (5.7,-1.3) rectangle (6.6,-1.3+\h);
  \node at (7.9,-1.3+0.5*\h) {piece $R$};
  \node at (9.8,-1.3+0.5*\h) {follower};
  \draw[decorate, decoration={brace, amplitude=3pt}] (5.3,-1.3+\h+0.06) -- (5.7,-1.3+\h+0.06)
    node[midway, above=3pt] {$\beta_d$};
  \draw[decorate, decoration={brace, mirror, amplitude=3pt}] (5.7,-1.36) -- (6.6,-1.36)
    node[midway, below=3pt] {deleted: $x_R$ blocks};
  \draw[densely dotted] (9.2,0) -- (9.2,-1.3+\h);
  \draw[densely dotted] (10.4,0) -- (10.4,-1.3+\h);
  \node[anchor=east] at (9.15,-0.48) {comparable followers};
  \draw[<->] (0.8,-0.48) -- node[below] {$\len P-\len R\ge\kappa_d$} (5.3,-0.48);
\end{scope}
%% after
\begin{scope}[yshift=-3.6cm]
  % P' = [0.8,5.3], junction at 1.2
  \draw (0,0) rectangle (7.1,\h);
  \node[left] at (0,0.5*\h) {$u_i$};
  \draw[fill=black!10] (0.8,0) rectangle (5.3,\h);
  \draw[line width=1.1pt, densely dashed] (1.2,-0.1) -- (1.2,\h+0.1);
  \node at (3.6,0.5*\h) {$P'$};
  \node at (6.2,0.5*\h) {follower};
  % R' = [2.3,5.3], junction at 2.7
  \draw (1.5,-1.15) rectangle (6.5,-1.15+\h);
  \node[left] at (1.5,-1.15+0.5*\h) {$u_j$};
  \draw[fill=black!10] (2.3,-1.15) rectangle (5.3,-1.15+\h);
  \draw[line width=1.1pt, densely dashed] (2.7,-1.25) -- (2.7,-1.15+\h+0.1);
  \node at (4.1,-1.15+0.5*\h) {$R'$};
  \node at (5.9,-1.15+0.5*\h) {follower};
  \draw[densely dotted] (5.3,0) -- (5.3,-1.15+\h);
  \draw[<->] (0.8,-0.42) -- node[below, xshift=-1mm] {$\ge\kappa_d$} (2.3,-0.42);
  \node[anchor=west, align=left] at (7.5,-0.4) {both pieces keep at least $\phi_d$ letters;\\
    junctions (dashed): where the blocks\\ were removed};
\end{scope}
\draw[->, semithick] (9.6,-1.9) -- (9.6,-2.9) node[midway, right] {admissible deletion};
\end{tikzpicture}
\caption{An admissible deletion (schematic). Two pieces $P$ and $R$ of the
same class whose followers are comparable can be aligned by a comparison, so
they must be shortened consistently. Each loses a number of blocks of $m_d$
letters (darker), starting $\beta_d$ letters after its start, where the states
of the run have become periodic with period $m_d$. Here $P$ is longer than
$R$ by at least the threshold $\kappa_d$; it loses more blocks, but remains longer
by at least $\kappa_d$. Pieces whose lengths differ by less than $\kappa_d$
lose the same number of blocks.}
\label{fig:compression}
\end{figure}


Inside a piece of period $d$ the states of the run are periodic with period
$m_d=d\cdot N!$ after $\alpha_d=Nd$ letters. So a block of $m_d$ letters
can be deleted from the interior of a piece without changing the later
states. Deleting blocks from one piece only would be visible to
comparisons that align this piece with a piece of the same class in another
word, or in the same word. Two pieces can be aligned by a comparison that
continues beyond them only if their followers are comparable, which leads to
the following conditions (\cref{fig:compression}). They use three lengths:
\[
  \beta_d=\alpha_d+m_d,\qquad\kappa_d=2\beta_d+m_d,\qquad
  \phi_d=6\beta_d+2m_d.
\]
The deleted blocks begin $\beta_d$ letters after the start of a piece: this
skips the transient $\alpha_d$ of the states and one more block. The number
$\kappa_d$ is the threshold for the difference of the lengths of two pieces
with comparable followers. Pieces that differ by less are shortened alike.
A difference of at least $\kappa_d$ is kept, and it ensures that where the
shorter piece begins, the states of the longer one are already periodic, so
that it does not matter how many blocks each of them has lost. Finally,
$\phi_d$ is the length that a shortened piece keeps; it leaves room for the
$\beta_d$ letters before the deleted blocks and for $5\beta_d+2m_d$ letters
after them, so that both ends of the piece and the letters next to them are
not affected. For instance, for $N=3$ and $d=2$ we get $\alpha_d=6$,
$m_d=12$, $\beta_d=18$, $\kappa_d=48$ and $\phi_d=132$.

\begin{definition}[Admissible deletion]\label{def:admissible}
An \emph{admissible deletion} of the class of $\zeta$, with $d=\len\zeta$,
deletes $x_P$ blocks of $m_d$ letters from every piece $P$ of the class,
starting $\beta_d$ letters after the start of $P$, such that
\begin{itemize}
\item a piece that is shortened keeps at least $\phi_d$ letters;
\item two pieces with comparable followers whose lengths differ by less
  than $\kappa_d$ lose the same number of blocks;
\item if $P$ and $R$ have comparable followers and $R$ is longer than $P$
  by at least $\kappa_d$, it remains longer by at least $\kappa_d$.
\end{itemize}
\end{definition}

\begin{lemma}[Admissible deletions]\label{lem:admissible}
The result of an admissible deletion dominates the original tuple. If the
tuple is a hole, the result is a hole with the same endpoints, its first
rejected cuts are the images of the original ones, and pieces of other
classes with periods at most $d$ and lengths at least $2d$ are not affected.
\end{lemma}

The numbers $x_P$ are the solution of another system of difference
constraints, and the least solution is small unless long chains of
constraints force a piece to stay long. Such a chain is a path in the graph
that joins pieces with comparable followers. Comparability of followers is
comparability in the forest of the prefix order, and a simple path in the
comparability graph of a forest of height $h$ has at most $2^h-1$ vertices.
For the pieces that end at or before the first rejected cut of their word,
or lie in words with mortal endpoints, comparable followers begin at cuts
with different states (\cref{lem:distinct}), so their height is at most
$k(N-1)$. The other pieces end in the rejected tails of immortal
coordinates, and their number is bounded by \cref{prop:run-cover}.
Splitting a path at these pieces gives a bound that is \emph{linear} in
their number.

\begin{lemma}[The cap]\label{lem:cap}
Suppose that every set of pieces of the class that have pairwise comparable
followers and are not among $r$ exceptional pieces has at most $h$ elements.
Then there is an admissible deletion after which every piece of the class
has length at most
\[
  \phi_d+m_d+\bigl((r+1)\,2^{h}-1\bigr)(\kappa_d+m_d).
\]
\end{lemma}

Starting from the hole of \cref{prop:run-cover}, we call the at most
$k(\Qc+3)$ pieces that cover its immortal tails the \emph{designated}
pieces. We compress the class of a designated piece of the largest period,
count the positions of the shortened designated pieces of this class among
the uncovered positions, and repeat with the remaining designated pieces.
After at most $k(\Qc+3)$ passes all immortal tails are short.

\begin{theorem}[Rejected tails, with the bound]\label{thm:tails-bound}
For $k\ge1$ and $N\ge1$, let $E_0$ and $D$ be as in \cref{prop:run-cover},
let $g=k(\Qc+3)$, and define
\[
  E_{r+1}=E_r+(g-r)\,\Bigl(\phi_D+m_D+\bigl((2(g-r)+E_r+1)\,2^{k(N-1)}-1\bigr)
  (\kappa_D+m_D)\Bigr).
\]
Then \cref{thm:tails} holds with
$\TailB(k,N)=E_g\le\expo_2\bigl((k+N+20)^{2k+5}\bigr)$ for $k\ge1$ and
$N\ge1$, and with $\TailB(0,N)=0$.
\end{theorem}

The closed form is proved in \cref{app:complexity-closed}.
